\documentclass[11pt]{article}

\usepackage[margin=1in]{geometry}
\usepackage[T1]{fontenc}
\usepackage[utf8]{inputenc}
\usepackage{lmodern}
\usepackage{microtype}
\usepackage{amsmath,amssymb}
\usepackage{booktabs}
\usepackage{tabularx}
\usepackage{array}
\usepackage{enumitem}
\usepackage{natbib}
\usepackage{xurl}
\usepackage[hidelinks]{hyperref}

\title{When Is a Formal Difference Admissible Evidence for an Individuation Claim?\\
\large Representation, Evidential Bridges, and Restricted Tests}
\author{Anonymous working draft}
\date{}

\newcommand{\Indist}{\mathrel{\equiv}}
\newcommand{\Dist}{\mathrel{\#}}
\newcolumntype{Y}{>{\raggedright\arraybackslash}X}

\begin{document}
\maketitle

\begin{abstract}
A formal difference between representations is not by itself evidence that represented targets are distinct. We formulate an individuation-specific admissibility framework separating representation features, target reference, evidential bridges, and admitted tests. A representation-level feature may contribute only when an independently justified experimental, causal, measurement, or semantic bridge connects it to the target; declaration alone does not supply warrant. A UMI sequencing case contrasts arbitrary read identifiers with experimentally grounded molecular tags. A small Lean artifact audits bridge dependence, relabeling invariance, and test-relative refinement. The framework states necessary, not sufficient, conditions for operational individuation claims.
\end{abstract}

\section{Introduction}

A formal model can contain two different representations without thereby providing evidence for two different target-level individuals. Variables, constructors, row identifiers, memory addresses, labels, and carrier elements can differ because of how a model is encoded rather than because the represented world contains two individuals.

The problem addressed here is not a general metaphysics of identity. It is an evidential question: when may a difference inside a representation be used as evidence for an operational individuation claim about the represented target?

Representation-level inequality alone is insufficient. A representation-level feature can contribute only when a defensible evidential bridge connects that feature to the target of the claim, and the resulting target-relevant difference is detectable under the tests admitted by the analysis.

The condition is necessary, not sufficient. A target-linked feature can still be unreliable, confounded, noisy, or badly calibrated. The framework asks a prior dependency question: has the argument shown why this formal difference is evidence about the target at all?

This yields a practical audit:
\begin{quote}
\textbf{If a representation-level identifier were arbitrarily reassigned while all independently justified target-relevant relations were held fixed, should the individuation conclusion change?}
\end{quote}
If the answer is yes, the argument owes an account of the evidential bridge from that identifier to the target.

The closest conceptual comparison is Nguyen's target-directed account of scientific representation and theoretical equivalence \citep{Nguyen2017RepresentationEquivalence}. The present paper asks a narrower downstream question: when a difference in the representation itself is offered as evidence for an individuation claim, what must connect that difference to the target before the inference is admissible?

\subsection{Contributions}

The paper makes three limited contributions.
\begin{enumerate}[leftmargin=*]
\item \textbf{An explicit evidential-bridge condition.} Representation-level features are separated from target-level properties, and the dependency connecting them is made explicit. Declaration exposes the dependency; scientific warrant must come from the relevant practice.
\item \textbf{A relabeling-invariance diagnostic with a scientific case.} A UMI sequencing case distinguishes arbitrary read identifiers from experimentally grounded molecular tags.
\item \textbf{A mechanized dependency audit.} A small Lean development checks bridge dependence, test-family refinement, and negative controls for representation-sensitive and semantically inert identity-like features.
\end{enumerate}

\section{Formal Setting}

\subsection{Representations and targets}

Let \(P\) be a set of representations, \(T\) a target domain, and
\[
r:P\to T
\]
a representation-to-target map. Distinct elements of \(P\) are not assumed to represent distinct elements of \(T\).

The map \(r\) is a semantic input. The paper does not infer \(r\) from observations. The formal core is restricted to one target value per representation; many-to-many and probabilistic correspondence are extensions outside the current scope.

\subsection{Representation features and evidential bridges}

Let
\[
f:P\to V
\]
be a representation-level feature and
\[
\phi:T\to V
\]
a target-level property. In the exact setting, write
\[
W(f,\phi)
\quad\Longleftrightarrow\quad
\forall p\in P,\; f(p)=\phi(r(p)).
\]

Then
\[
W(f,\phi)\land f(a)\neq f(b)
\quad\Longrightarrow\quad
r(a)\neq r(b).
\]

The implication is elementary. Its value is structural: the bridge makes explicit what must connect a presentation-level difference to the target. Declaration does not create epistemic warrant. In a scientific application, accepting the bridge requires an independent experimental, causal, measurement, calibration, or semantic justification.

\subsection{Admitted tests}

Let \(Q\) be a class of target-sensitive tests and \(O(q,t)\) the outcome of test \(q\) on target \(t\). Let \(A\subseteq Q\) be a family admitted by an analysis. Admission itself requires an independently defensible connection between the test and the target claim.

Define
\[
a\Indist_A b
\quad\Longleftrightarrow\quad
\forall q\in A,\;O(q,r(a))=O(q,r(b)),
\]
and
\[
a\Dist_A b
\quad\Longleftrightarrow\quad
\exists q\in A:\;O(q,r(a))\neq O(q,r(b)).
\]

These relations are operational and test-relative. They are not definitions of metaphysical numerical identity.

\section{Scientific Case: UMI-Based Molecular Counting}

PCR amplification can produce many downstream sequencing reads from one pre-amplification DNA or RNA molecule. A hypothetical data set containing 100 read records therefore does not, merely by containing 100 records, support the conclusion that 100 source molecules were present.

\citet{Kivioja2012UMI} introduced unique molecular identifiers (UMIs) to support molecule counting by attaching molecular tags before amplification. \citet{SmithHegerSudbery2017UMITools} show why the tags still require error-aware analysis: UMI sequences can themselves contain sequencing errors, and naive token equality can misidentify PCR duplicates.

An ordinary software read identifier distinguishes read records. Its inequality is representation-level and, by itself, supplies no bridge to source-molecule plurality.

A pre-amplification UMI differs because the experimental procedure is designed to connect the tag to a source molecule before PCR generates downstream multiplicity. In the idealized model, \(P\) is the set of read records, \(T\) the set of pre-amplification molecules, \(r(p)\) the source molecule of read \(p\), and admitted evidence includes the molecular tag together with alignment and contextual information.

Actual investigators do not directly observe \(r\). Source assignment is inferred through the protocol, alignment, barcode statistics, and error assumptions. The framework therefore audits a proposed source-assignment and evidential model; it does not reconstruct deduplication from raw data.

Target relevance is also not infallibility. Barcode collisions, sequencing errors, and other failures mean that same-UMI and different-UMI observations cannot be treated as unrestricted numerical-identity rules. Experimental provenance makes the tag evidentially relevant; the error model determines how particular observations should be interpreted.

\section{Test-Relative Distinguishability}

Let \(A_c\subseteq A_f\). Under fixed exact outcome semantics,
\[
a\Indist_{A_f}b\Longrightarrow a\Indist_{A_c}b.
\]
The converse need not hold.

The interpretation is limited:
\begin{quote}
\textbf{Under fixed exact semantics, enlarging a justified test family can refine the induced operational classification.}
\end{quote}
Noise, statistical updating, model revision, and measurement disturbance are outside this exact result.

\section{Mechanized Dependency Audit}

The Lean artifact is a dependency audit, not a claim of deep mathematical novelty.

The new bridge layer represents a presentation-level Boolean feature together with a referent-level Boolean property through which that feature factors. The development proves that equal grounding forces equality of a bridged feature and that a difference in a bridged feature entails a grounding difference. It also proves that the deliberately representation-sensitive encoding probe has no such bridge, because it separates two presentations with equal grounding, while a positive-control ground-tracking feature does have one.

The older controls remain as supporting checks. Same-reference encodings preserve their target-sensitive profile. A restricted test family can fail to separate a pair that a richer family separates. An arbitrary identity-like token carrier is semantically inert unless it enters the declared target semantics.

The results are simple by design. Their purpose is to expose hidden dependencies rather than to replace substantive scientific justification.

\section{Relation to Existing Work}

\subsection{Scientific representation and theoretical equivalence}

\citet{Suarez2004Inferential} and \citet{Contessa2007Surrogative} already treat scientific representation as target-directed and inferential. \citet{Nguyen2017RepresentationEquivalence} is the closest comparison: model and theoretical equivalence can be analyzed by whether models license the same claims about the same targets.

The present paper claims no novelty for target-directed inference or same-target claim licensing. Its narrower question begins when a formal difference inside a representation is itself offered as evidence for target plurality. The bridge condition isolates that dependency: the feature must have a justified route to a target-level property rather than merely be formally discriminating.

\subsection{Individuation in scientific practice}

Practice-oriented work asks how scientists count, track, separate, manipulate, and present entities in concrete settings \citep{BuenoChenFagan2018,Waters2018AskNot,Love2018ExperimentalPractice}. \citet{Chen2018ExperimentalIndividuation} distinguishes ontological and epistemological modes of experimental individuation and explicitly discusses presentation.

The present contribution is narrower: when a representational device is used in an individuation practice, what justifies treating differences in that device as evidence about the target?

\subsection{Technical neighbors}

Representation independence \citep{Mitchell1986RepresentationIndependence}, observational and behavioral equivalence \citep{HennessyMilner1985,Rutten2000UniversalCoalgebra}, and work on identity and discernibility \citep{LadymanLinneboPettigrew2012,DieksVersteegh2008} constrain the interpretation of the present results. No new general equivalence theory or solution to numerical identity is claimed.

Nor does the negative result imply that representation-only or surplus structure is useless \citep{NguyenTehWells2020Surplus}. The claim is evidential, not eliminativist.

\section{Scope and Limitations}

The framework leaves several problems open.

First, \(T\) and \(r\) are inputs. It does not derive target ontology or solve uncertain source assignment.

Second, the bridge condition is exact and deterministic. Real scientific warrants can be probabilistic, error-prone, and model-dependent.

Third, target relevance and test discrimination are necessary but not sufficient for epistemic warrant. Reliability, calibration, confounding, background assumptions, and statistical decision rules remain substantive scientific questions.

Fourth, the formal core uses a function \(r:P\to T\). Many-to-many, distributed, and probabilistic representation-target relations are not treated here.

Fifth, observation and intervention are not equated. Different test classes can induce different operational partitions.

Finally, operational distinguishability is not metaphysical numerical identity, and synchronic discrimination does not establish diachronic persistence.

\section{Relabeling Invariance as a Practical Audit}

For an identifier-like feature used in an individuation argument, ask whether arbitrary reassignment would alter the conclusion while independently justified target relations were held fixed. If it would, ask what experimental, causal, measurement, or semantic bridge explains why.

A representation-only discriminator changes bookkeeping or encoding while leaving the target-relevant model fixed. A bridged discriminator participates in a justified process connecting representation to target. It may therefore contribute evidence, subject to the reliability and error assumptions of that process.

The UMI case instantiates the contrast. Read identifiers fail the molecular relabeling audit. Pre-amplification UMIs are designed to pass it because their evidential role depends on the physical tagging history, not on identifierhood.

\section{Claim-to-Artifact Map}

The public development contains a broader set of elementary checks than the argument needs in the main text.

\begin{center}
\small
\begin{tabularx}{\linewidth}{@{}>{\raggedright\arraybackslash}p{0.45\linewidth}Y@{}}
\toprule
Claim family & Public Lean checks \\
\midrule
Representation-sensitive discrimination versus same-reference invariance &
\path{encodingProbe_separates_sameGrounding};
\path{sameGrounding_sameGroundedObservation};
\path{sameGrounding_sameGroundedProfile};
\path{sameReferent_encodings_groundedlyIndistinguishable} \\
Test-family refinement &
\path{groundedIndist_of_access_inclusion};
\path{coarse_indistinguishable};
\path{fine_distinguishable} \\
Target-sensitive positive control &
\path{groundedOutcomeDifference_impliesGroundDifference};
\path{identityProbe_groundedDifference};
\path{derivedGroundDifference} \\
Token/access negative controls &
\path{tokenizedGroundedIndist_iff_base};
\path{identityToken_variation_preserves_groundedClassification};
\path{differentIdentityTokens_sameGroundedClassification};
\path{identityTokenCannotMask_groundedDifference};
\path{decorativeInterfaceTag_irrelevant} \\
Evidential-bridge dependency &
\path{bridgedFeature_sameGrounding_sameValue};
\path{bridgedFeatureDifference_impliesGroundDifference};
\path{encodingProbe_hasNoGroundingBridge};
\path{groundFeature_hasGroundingBridge} \\
\bottomrule
\end{tabularx}
\end{center}

\section{Conclusion}

A formal difference is not admissible evidence for an individuation claim merely because it distinguishes two representations.

The missing step is an evidential bridge. A representation-level feature may contribute to target-level individuation only when an independently justified experimental, causal, measurement, or semantic relation connects that feature to the target. A difference must then survive the relevant admitted tests. These are necessary conditions, not a complete epistemology of individuation.

UMI-based molecular counting shows why this matters. A read identifier and a molecular barcode can have the same formal type while playing very different evidential roles. The difference is explained by provenance and error-aware measurement practice, not by the syntax of an identifier.

The formal contribution is correspondingly modest: make the bridge explicit, make relabeling invariance testable, and make hidden dependence on representation-level identity mechanically auditable.

\bibliographystyle{plainnat}
\bibliography{references}

\end{document}
